% !TeX root = ../main.tex
\section{Gruppi: assiomi e prime proprietà}
% A–L L17 pp. 3–6 e L18 p. 2; M–Z L17 pp. 1–3.
Un anello possiede due operazioni; un gruppo ne considera una sola. La struttura dipende sia dall'insieme sia dall'operazione scelta: gli interi sono un gruppo per la somma, ma non per il prodotto.

\begin{definizione}
Un \textbf{gruppo} è una coppia $(G,\cdot)$ con $G\ne\varnothing$ e un'operazione binaria $\cdot:G\times G\to G$ tale che:
\begin{enumerate}
\item $(ab)c=a(bc)$ per ogni $a,b,c\in G$;
\item esiste un \textbf{elemento neutro} $1_G\in G$ con $1_Ga=a1_G=a$ per ogni $a\in G$;
\item per ogni $a\in G$ esiste un \textbf{inverso} $b\in G$ con $ab=ba=1_G$.
\end{enumerate}
Il gruppo è \textbf{abeliano}, o \textbf{commutativo}, se inoltre $ab=ba$ per ogni $a,b\in G$. Il prodotto $ab$ abbrevia $a\cdot b$; la commutatività non è richiesta negli assiomi di gruppo.
\end{definizione}

\begin{proposizione}\label{prop:gruppo-prime}
In un gruppo valgono:
\begin{enumerate}[label=\arabic*.]
\item Il neutro e l\textquotesingle inverso di ciascun elemento sono unici. Indicato l\textquotesingle inverso di $a$ con $a^{-1}$, si ha $(a^{-1})^{-1}=a$.
\item $ac=bc\Rightarrow a=b$ e $ca=cb\Rightarrow a=b$ (cancellazione).
\item $(ab)^{-1}=b^{-1}a^{-1}$.
\end{enumerate}
\end{proposizione}
\begin{proof}
\begin{enumerate}[label=\arabic*.]
\item Se $e,e'$ sono neutri, $e=ee'=e'$. Se $b,c$ sono inversi di $a$, allora $b=b(ac)=(ba)c=c$. Le identità $aa^{-1}=a^{-1}a=1_G$ mostrano anche che $a$ è l'inverso di $a^{-1}$.
\item Da $ac=bc$ si ottiene $(ac)c^{-1}=(bc)c^{-1}$ e, per associatività, $a=b$. Da $ca=cb$ si moltiplica invece a sinistra per $c^{-1}$. Non si scambia l'ordine dei fattori.
\item Infine
\[
(ab)(b^{-1}a^{-1})=a(bb^{-1})a^{-1}=1_G,
\qquad (b^{-1}a^{-1})(ab)=b^{-1}(a^{-1}a)b=1_G.
\]
L'unicità dell'inverso dà la formula per $(ab)^{-1}$: occorre invertire anche l'ordine.
\end{enumerate}
\end{proof}

\subsection{Potenze e notazione additiva}
\begin{definizione}
Per $a\in G$ e $n\ge1$, $a^n$ è il prodotto di $n$ fattori uguali ad $a$. Si pone inoltre
\[
a^0=1_G,\qquad a^{-n}=(a^{-1})^n\quad(n\ge1).
\]
In un gruppo abeliano si può usare la \textbf{notazione additiva}: l'operazione si scrive $+$, il neutro $0$, l'inverso di $a$ si scrive $-a$, e la potenza $a^n$ diventa il multiplo $na$. Per esempio, in $(\Z,+)$, $3a=a+a+a$.
\end{definizione}

\begin{proposizione}\label{prop:potenze-gruppo}
Per ogni $a\in G$ e $m,n\in\Z$,
\begin{enumerate}[label=\arabic*.]
\item $a^ma^n=a^{m+n}$.
\item $(a^n)^{-1}=a^{-n}$.
\end{enumerate}
In particolare due potenze dello stesso elemento commutano, anche se $G$ non è abeliano.
\end{proposizione}
\begin{proof}
\begin{enumerate}[label=\arabic*.]
\item La concatenazione dei fattori dà la prima identità per $m,n\ge0$. Ripetendo la formula dell'inverso di un prodotto si ha $(a^k)^{-1}=(a^{-1})^k$ per $k\ge1$.

Se $m\ge0$ e $n=-k<0$, per $m\ge k$ si scrive $a^m=a^{m-k}a^k$ e si cancella $a^k$ con $(a^k)^{-1}$. Per $m<k$ si scrive $a^k=a^{k-m}a^m$, quindi $a^m(a^k)^{-1}=(a^{k-m})^{-1}=a^{m-k}$. Il caso $m<0,n\ge0$ si ottiene con le stesse cancellazioni sull'altro lato. Se entrambi sono negativi, si applica la concatenazione alle potenze di $a^{-1}$. La formula vale così per tutti gli interi.
\item Ponendo $m=-n$ si verifica l'identità sull'inverso; infine $m+n=n+m$ dà la commutazione.
\end{enumerate}
\end{proof}

\subsection{Gruppi ottenuti dagli anelli}
\begin{proposizione}
\begin{enumerate}[label=\arabic*.]
\item Se $R$ è un anello, $(R,+)$ è un gruppo abeliano.
\item Se $R$ ha unità, $(U(R),\cdot)$ è un gruppo, abeliano quando $R$ è commutativo.
\item Se $K$ è un campo, $U(K)=K\setminus\{0\}$.
\end{enumerate}
\end{proposizione}
\begin{proof}
\begin{enumerate}[label=\arabic*.]
\item Gli assiomi additivi di anello sono esattamente quelli di un gruppo abeliano.
\item Per il prodotto delle unità, l'operazione è interna perché, per $a,b\in U(R)$, $b^{-1}a^{-1}$ è un inverso bilatero di $ab$. L'associatività è ereditata da $R$, $1_R$ è invertibile e l'inverso di un'unità è ancora un'unità. Se $R$ è commutativo lo è anche il prodotto ristretto.
\item In un campo ogni elemento non nullo è invertibile, mentre lo zero non lo è perché $1_K\ne0_K$.
\end{enumerate}
\end{proof}

\begin{esempio}
$(\Z,+)$, $(\Z/25\Z,+)$ e $(U(\Z/25\Z),\cdot)$ sono gruppi abeliani. I due gruppi modulo $25$ hanno rispettivamente $25$ e $20$ elementi: cambiano sia l'insieme sia l'operazione. In generale $U(\Z/m\Z)$ ha $\varphi(m)$ elementi. Poiché $13$ è primo, $U(\Z/13\Z)$ è il gruppo moltiplicativo delle dodici classi non nulle.

$(\Z,\cdot)$, $(\Q,\cdot)$ e $(\R,\cdot)$ non sono gruppi: contengono lo zero, che non ha inverso moltiplicativo. Sono invece gruppi abeliani $(\Q\setminus\{0\},\cdot)$ e $(\R\setminus\{0\},\cdot)$.
\end{esempio}

\section{Permutazioni e trasformazioni affini}
\subsection{Il gruppo simmetrico}
% A–L L17 pp. 7–8; M–Z L17 pp. 3–5.
Le operazioni puntuali sulle funzioni e la composizione sono operazioni diverse. Nell'anello $\R^{\R}$ con prodotto puntuale, le unità sono le funzioni che non si annullano mai e formano un gruppo abeliano; per la composizione, invece, le funzioni invertibili sono esattamente le biezioni. Tutte le funzioni $\R\to\R$ con la sola composizione formano un monoide: la composizione è associativa e ha neutro $\operatorname{id}_{\R}$, ma non tutte le funzioni hanno inversa.

\begin{definizione}
Per un insieme non vuoto $X$, si indica con
\[
S_X=\operatorname{Sym}(X)=\operatorname{Perm}(X)
   =\{\sigma:X\to X\mid\sigma\text{ è biettiva}\}
\]
l'insieme delle \textbf{permutazioni} di $X$. Con la composizione forma il \textbf{gruppo simmetrico} su $X$. Se $X=\{1,\ldots,n\}$, si scrive $S_n$.
\end{definizione}

\begin{proposizione}\label{prop:simmetrico}
\begin{enumerate}[label=\arabic*.]
\item $S_X$ è un gruppo con neutro $\operatorname{id}_X$ e inversa la funzione inversa.
\item Per $n\ge1$, $|S_n|=n!$.
\item Se $X$ ha almeno tre elementi, $S_X$ non è abeliano.
\end{enumerate}
\end{proposizione}
\begin{proof}
\begin{enumerate}[label=\arabic*.]
\item La composta di due biezioni è biettiva; la composizione è associativa, l'identità è biettiva e l'inversa di una biezione è biettiva. Sono le proprietà già dimostrate per le funzioni.
\item Per scegliere una biezione di $\{1,\ldots,n\}$, l'immagine di $1$ ha $n$ possibilità, quella di $2$ ne ha $n-1$ perché deve essere diversa, e così via fino all'ultima scelta obbligata. Il conteggio è $n(n-1)\cdots1=n!$.
\item Scelti tre punti distinti $x,y,z$, siano $s$ lo scambio di $x,y$ e $t$ lo scambio di $y,z$, lasciando fissi gli altri punti. Allora $(s\circ t)(x)=y$, mentre $(t\circ s)(x)=z$: le due composte sono diverse. Per uno o due punti, invece, il gruppo è abeliano (rispettivamente il solo neutro, oppure neutro e unico scambio).
\end{enumerate}
\end{proof}

\Needspace{30\baselineskip}
\begin{esempio}[Leggere e comporre le permutazioni]
La notazione su due righe elenca sopra gli argomenti e sotto le immagini. In $S_4$, siano
\[
\sigma=\begin{pmatrix}1&2&3&4\\1&3&2&4\end{pmatrix},\qquad
\tau=\begin{pmatrix}1&2&3&4\\4&2&1&3\end{pmatrix}.
\]
Si ha $\sigma^2=\operatorname{id}$, quindi $\sigma^{-1}=\sigma$. Nella composta $\sigma\circ\tau$ si applica prima $\tau$, poi $\sigma$:
\[
\begin{array}{c|rrrr}
x&1&2&3&4\\\hline
\tau(x)&4&2&1&3\\
\sigma(\tau(x))&4&3&1&2
\end{array}
\]
Per seguire un singolo punto, il percorso è
\begin{center}
\begin{tikzpicture}[>=Stealth,baseline]
\node (a) at (0,0) {$2$}; \node (b) at (2.2,0) {$2$}; \node (c) at (4.4,0) {$3$};
\draw[->] (a)--node[above]{$\tau$}(b); \draw[->] (b)--node[above]{$\sigma$}(c);
\end{tikzpicture}
\end{center}
Invertendo l'ordine si ottiene invece
\[
\sigma\circ\tau=\begin{pmatrix}1&2&3&4\\4&3&1&2\end{pmatrix},\qquad
\tau\circ\sigma=\begin{pmatrix}1&2&3&4\\4&1&2&3\end{pmatrix}.
\]
Un secondo esempio, con $\rho=\left(\begin{smallmatrix}1&2&3&4\\3&2&4&1\end{smallmatrix}\right)$, dà
\[
\rho\circ\sigma=\begin{pmatrix}1&2&3&4\\3&4&2&1\end{pmatrix},\qquad
\sigma\circ\rho=\begin{pmatrix}1&2&3&4\\2&3&4&1\end{pmatrix}.
\]
Anche qui l'ordine delle applicazioni modifica il risultato.
\end{esempio}

\subsection{Il gruppo affine reale}
% A–L L17 pp. 9–10 e L18 p. 1 (pagina ripetuta); M–Z L17 pp. 5–6.
\begin{definizione}
Per $a\in\R\setminus\{0\}$ e $b\in\R$, poniamo $f_{a,b}(x)=ax+b$. L'insieme
\[
\operatorname{Aff}(\R)=\{f_{a,b}:\R\to\R\mid a\ne0,\ b\in\R\}
\]
è l'insieme delle \textbf{trasformazioni affini reali invertibili}; lo consideriamo con la composizione. I coefficienti sono univoci: $b=f_{a,b}(0)$ e $a=f_{a,b}(1)-f_{a,b}(0)$.
\end{definizione}

\begin{proposizione}\label{prop:affini}
\begin{enumerate}[label=\arabic*.]
\item $\operatorname{Aff}(\R)$ è un gruppo non abeliano e
\[
f_{a,b}\circ f_{a',b'}=f_{aa',\,ab'+b},\qquad
1_{\operatorname{Aff}(\R)}=f_{1,0},\qquad
f_{a,b}^{-1}=f_{a^{-1},\,-ba^{-1}}.
\]
\item Due trasformazioni $f_{a,b}$ e $f_{a',b'}$ commutano se e solo se
\[
(1-a')b=(1-a)b'.
\]
\end{enumerate}
\end{proposizione}
\begin{proof}
\begin{enumerate}[label=\arabic*.]
\item Applicando prima $f_{a',b'}$ e poi $f_{a,b}$ si ottiene $a(a'x+b')+b=aa'x+ab'+b$. Poiché $aa'\ne0$, la composizione appartiene ancora all'insieme. L'associatività è quella della composizione di funzioni e $f_{1,0}$ è l'identità. Risolvendo $y=ax+b$ si trova $x=a^{-1}y-ba^{-1}$; sostituendo, questa funzione è inversa su entrambi i lati. Quindi ogni elemento è una biezione e gli assiomi di gruppo sono verificati.
\item Le due composte hanno lo stesso coefficiente $aa'=a'a$; coincidono se e solo se $ab'+b=a'b+b'$, equivalente al criterio enunciato. Non vale per ogni coppia: $f_{2,1}\circ f_{1,2}$ manda $x$ in $2x+5$, mentre $f_{1,2}\circ f_{2,1}$ lo manda in $2x+3$.
\end{enumerate}
\end{proof}
